\documentclass[11pt]{article}
\usepackage[margin=1in]{geometry}
\usepackage{amsmath,amssymb,amsthm}
\usepackage{graphicx}
\usepackage{hyperref}

\newtheorem{theorem}{Theorem}
\newtheorem{lemma}{Lemma}
\newtheorem{corollary}{Corollary}
\newtheorem{remark}{Remark}

\newcommand{\E}{\mathbb E}
\newcommand{\Pp}{\mathbb P}
\newcommand{\R}{\mathbb R}
\newcommand{\cB}{\mathcal B}
\newcommand{\cL}{\mathcal L}
\newcommand{\norm}[1]{\left\lVert #1\right\rVert}
\newcommand{\simequiv}{\simeq}

\title{Discrete Aggregate Domination in \(L^q\)-Spaces}
\author{Agent \texttt{agent\_lane\_19}}
\date{2026-06-16}

\begin{document}
\maketitle

\section*{Status}

This packet gives a candidate partial result, likely valid, for Remark 12.10 of
Yaroslavtsev's paper \emph{Local characteristics and tangency of vector-valued
martingales} \cite{Yaroslavtsev2019}.  It proves the discrete
independent-increment aggregate-domination comparison in every \(L^q\)-space
for \(1<p,q<\infty\).

It is partial relative to the full source problem because it does not yet
package the full general-local-martingale extension, and it does not cover the
endpoint \(p=1\).

\section*{Source Problem}

The source is arXiv:1907.11588, PDF page 105, Remark 12.10.  The source crop
below states that the general-local-martingale case is open and reduces the
missing accessible-jump issue to a discrete independent-increment comparison.

\begin{center}
\includegraphics[width=\textwidth]{figures/open_problem_crop.png}
\end{center}

The next page adds that, by symmetrization, one may assume the variables are
symmetric, and that even this symmetric discrete case was not known to the
author.

\begin{center}

\includegraphics[width=\textwidth]{figures/open_problem_continuation_crop.png}
\end{center}

\section*{Main Result}

\begin{theorem}\label{thm:Lq-discrete}
Let \(1<p,q<\infty\), let \((S,\Sigma,\rho)\) be a measure space, and put
\(X=L^q(S)\).  Let \((\xi_i)_{i=1}^n\) and \((\eta_i)_{i=1}^m\) be finite
families of independent mean-zero \(X\)-valued random variables.  Assume that
\[
   \sum_i \Pp(\xi_i\in B)
   \leq
   \sum_j \Pp(\eta_j\in B)
   \qquad
   (B\in\cB(X),\;0\notin B).
\]
Then
\[
   \E\norm{\sum_i\xi_i}_{L^q(S)}^p
   \leq
   C_{p,q}\,
   \E\norm{\sum_j\eta_j}_{L^q(S)}^p .
\]
\end{theorem}

\begin{corollary}\label{cor:maximal}
Under the assumptions of Theorem \ref{thm:Lq-discrete},
\[
   \E\max_{1\leq k\leq n}\norm{\sum_{i=1}^k\xi_i}_{L^q(S)}^p
   \leq
   C_{p,q}\,
   \E\max_{1\leq k\leq m}\norm{\sum_{j=1}^k\eta_j}_{L^q(S)}^p .
\]
\end{corollary}

\section*{Idea of the Proof}

The previous \(\ell_q\) packet worked at \(p=q\) because the \(q\)-moment of
the lattice square function is additive across coordinates.  Off the natural
exponent, that additivity disappears.  The correct replacement is the sharp
\(L^q\)-valued Rosenthal theorem of Dirksen and Yaroslavtsev
\cite{DirksenYaroslavtsev2017}.

For independent mean-zero \(L^q\)-valued variables \(z_i\), their theorem says
that
\[
   \left(\E\norm{\sum_i z_i}_{L^q}^p\right)^{1/p}
   \simequiv_{p,q}
   \norm{(z_i)}_{s_{p,q}},
\]
where \(s_{p,q}\) is built, by sums and intersections, from three gauges:
\[
  S_q(z)=\left\|\left(\sum_i\E |z_i|^2\right)^{1/2}\right\|_{L^q},
  \quad
  D_{q,q}(z)=\left(\sum_i\E\norm{z_i}_{L^q}^q\right)^{1/q},
  \quad
  D_{p,q}(z)=\left(\sum_i\E\norm{z_i}_{L^q}^p\right)^{1/p}.
\]
These gauges depend only on the aggregate occurrence measure
\(\sum_i\mathcal L(z_i)\).

The only point that is not completely formal is the presence of sum spaces,
such as \(S_q+D_{q,q}\).  A near-optimal decomposition of the dominating
variables induces a measure on split triples \((x,y_1,\ldots,y_r)\), where
\(x=\sum_\ell y_\ell\).  If the aggregate law of \((\xi_i)\) is dominated by
that of \((\eta_j)\), we restrict this split measure by the Radon-Nikodym
derivative of the smaller aggregate law.  This transports the decomposition to
the \(\xi_i\)'s and decreases every basic gauge.  Thus the full Rosenthal norm
is monotone under aggregate domination.

\section*{The Rosenthal Norms}

Let \(X=L^q(S)\), \(1<p,q<\infty\).  For a finite sequence
\(z=(z_i)\) of \(X\)-valued random variables define
\[
  S_q(z)=\left\|\left(\sum_i\E |z_i|^2\right)^{1/2}\right\|_{L^q(S)},
\]
\[
  D_{r,q}(z)=\left(\sum_i\E\norm{z_i}_{L^q(S)}^r\right)^{1/r},
  \qquad r\in\{p,q\}.
\]
The Dirksen-Yaroslavtsev independent Rosenthal theorem
\cite[Theorem 2.1]{DirksenYaroslavtsev2017} says that for independent
mean-zero \(z_i\),
\[
   \left(\E\norm{\sum_i z_i}_{L^q(S)}^p\right)^{1/p}
   \simequiv_{p,q}
   \norm{z}_{s_{p,q}},
\]
where \(s_{p,q}\) is one of the following six sum/intersection spaces:
\[
\begin{array}{ll}
S_q\cap D_{q,q}\cap D_{p,q}, & 2\le q\le p<\infty,\\
S_q\cap(D_{q,q}+D_{p,q}), & 2\le p\le q<\infty,\\
(S_q\cap D_{q,q})+D_{p,q}, & 1<p<2\le q<\infty,\\
(S_q+D_{q,q})\cap D_{p,q}, & 1<q<2\le p<\infty,\\
S_q+(D_{q,q}\cap D_{p,q}), & 1<q\le p\le2,\\
S_q+D_{q,q}+D_{p,q}, & 1<p\le q\le2.
\end{array}
\]
Intersections carry the maximum norm and sums carry the usual infimum norm.

\section*{Aggregate Monotonicity}

For a finite sequence \(z=(z_i)\), let
\[
   \mu_z=\sum_i\mathcal L(z_i)|_{X\setminus\{0\}}
\]
be its aggregate occurrence measure.  The next lemma is the key point.

\begin{lemma}\label{lem:monotone}
Let \(a=(a_i)\) and \(b=(b_j)\) be finite sequences of \(X\)-valued random
variables with \(\mu_a\leq\mu_b\).  Then
\[
   \norm{a}_{s_{p,q}}\leq \norm{b}_{s_{p,q}} .
\]
\end{lemma}

\begin{proof}
First consider the three basic gauges.  Since \(\mu_a\leq\mu_b\),
\[
  \int_X |x|^2\,d\mu_a(x)
  \leq
  \int_X |x|^2\,d\mu_b(x)
\]
pointwise on \(S\), and therefore \(S_q(a)\leq S_q(b)\).  Similarly,
\[
  D_{r,q}(a)^r
  =
  \int_X\norm{x}_{L^q}^r\,d\mu_a(x)
  \leq
  \int_X\norm{x}_{L^q}^r\,d\mu_b(x)
  =
  D_{r,q}(b)^r
\]
for \(r=p,q\).

This proves the assertion for every expression involving only intersections.
It remains to handle sums.  We prove the required monotonicity recursively for
any expression built from the three basic gauges using \(\cap\) and \(+\).
The intersection step is immediate from the maximum norm.  For the sum step,
write the expression as \(E+F\).  Fix \(\varepsilon>0\), and choose a
decomposition \(b=u+v\) such that
\[
   \norm{u}_{E}+\norm{v}_{F}
   \leq
   \norm{b}_{E+F}+\varepsilon .
\]
Here \(u=(u_j)\), \(v=(v_j)\), and \(b_j=u_j+v_j\).

The decomposition induces a finite measure \(\lambda\) on
\[
   \{(x,y,z)\in X^3:\ x=y+z\}
\]
by
\[
   \lambda=\sum_j\mathcal L(b_j,u_j,v_j)|_{X\setminus\{0\}\times X\times X}.
\]
Its \(x\)-marginal is \(\mu_b\).  Let \(h=d\mu_a/d\mu_b\), \(0\leq h\leq1\),
and replace \(\lambda\) by \(h(x)\lambda(dx,dy,dz)\).  Disintegrating this
restricted measure over its \(x\)-marginal \(\mu_a\), and enlarging the
probability spaces of the \(a_i\)'s if necessary, we obtain random variables
\(a_i^E,a_i^F\) such that
\[
   a_i=a_i^E+a_i^F
\]
and the aggregate occurrence measures of \((a_i^E)\) and \((a_i^F)\) are
dominated by those of \(u\) and \(v\), respectively.  By the induction
hypothesis,
\[
   \norm{(a_i^E)}_E\leq\norm{u}_E,
   \qquad
   \norm{(a_i^F)}_F\leq\norm{v}_F .
\]
Thus
\[
   \norm{a}_{E+F}
   \leq
   \norm{(a_i^E)}_E+\norm{(a_i^F)}_F
   \leq
   \norm{b}_{E+F}+\varepsilon .
\]
Letting \(\varepsilon\downarrow0\) completes the recursive proof.

The above disintegration is first applied on the separable closed subspace
generated by the essentially separable ranges of the finite family involved;
standard approximation then removes boundedness or simplicity restrictions.
\end{proof}

\section*{Proof of the Main Result}

Let \(\mu_\xi=\sum_i\mathcal L(\xi_i)|_{X\setminus\{0\}}\) and
\(\mu_\eta=\sum_j\mathcal L(\eta_j)|_{X\setminus\{0\}}\).  The hypothesis is
exactly \(\mu_\xi\leq\mu_\eta\).  By Lemma \ref{lem:monotone},
\[
   \norm{(\xi_i)}_{s_{p,q}}
   \leq
   \norm{(\eta_j)}_{s_{p,q}} .
\]
The independent \(L^q\)-valued Rosenthal theorem gives
\[
\begin{aligned}
  \left(\E\norm{\sum_i\xi_i}_{L^q}^p\right)^{1/p}
  &\lesssim_{p,q}
  \norm{(\xi_i)}_{s_{p,q}} \\
  &\leq
  \norm{(\eta_j)}_{s_{p,q}} \\
  &\lesssim_{p,q}
  \left(\E\norm{\sum_j\eta_j}_{L^q}^p\right)^{1/p}.
\end{aligned}
\]
Raising to the \(p\)-th power proves Theorem \ref{thm:Lq-discrete}.

For Corollary \ref{cor:maximal}, apply Doob's \(L^p\) maximal inequality to
the martingale \(\sum_{i\leq k}\xi_i\), use Theorem \ref{thm:Lq-discrete}, and
then bound the terminal \(\eta\)-sum by the \(\eta\)-maximum.

\section*{Verification and Limitations}

The proof uses the lower and upper bounds in the Dirksen-Yaroslavtsev
Rosenthal theorem, so it is restricted to \(1<p<\infty\).  Their paper records
some endpoint upper bounds, but the two-sided norm identification needed here
is not available in this packet at \(p=1\).

This packet proves the discrete independent-increment bottleneck from Remark
12.10 in \(L^q\)-spaces.  A full general-local-martingale \(L^q\) theorem
should now be within reach, but it requires a separate audit combining this
jump comparison with continuous Gaussian covariance domination and the source
paper's decoupled-tangent reduction.  The present packet deliberately does not
claim that final step.

The code file \texttt{code/grouping\_simulation.py} is only a finite sanity
check.  It compares several groupings of the same aggregate coordinate jump
law in \(\ell_q^N\).  The tested ratios remain dimension-stable, supporting
the monotonicity proof but not serving as evidence in place of it.

\section*{Novelty Check}

Cheap run indexes contained the previous \(c_0\) counterexample, Hilbert
positive packets, and the \(\ell_q\) natural-exponent packet, but no
\(L^q\) all-\(p>1\) aggregate-domination packet for arXiv:1907.11588.  Web
searches for exact phrases around \texttt{characteristically dominated},
\(\ell_q\), \(L^q\), and \texttt{aggregate domination} did not locate an
existing statement of this consequence.  The decisive external input is the
Dirksen-Yaroslavtsev \(L^q\)-valued Burkholder-Rosenthal theorem
\cite{DirksenYaroslavtsev2017}; the monotonicity lemma above is the additional
identification that connects that theorem to Remark 12.10.

\section*{Human Review Recommendation}

Review Lemma \ref{lem:monotone} first, especially the transfer of
near-optimal sum-space decompositions by disintegration and
Radon-Nikodym restriction.  If that lemma is accepted, the rest of the proof is
an immediate application of the cited Rosenthal theorem.

The next target is to combine this packet with the full decoupled-tangent
reduction in \cite{Yaroslavtsev2019}.  The jump part should follow from the
present theorem; the continuous part should follow from Gaussian covariance
domination, but it should be audited separately before promoting a full
\(L^q\) general-local-martingale result.

\begin{thebibliography}{9}

\bibitem{Yaroslavtsev2019}
I. S. Yaroslavtsev,
\emph{Local characteristics and tangency of vector-valued martingales},
arXiv:1907.11588.

\bibitem{DirksenYaroslavtsev2017}
S. Dirksen and I. S. Yaroslavtsev,
\emph{\(L^q\)-valued Burkholder-Rosenthal inequalities and sharp estimates for
stochastic integrals},
arXiv:1707.00109.

\end{thebibliography}

\end{document}
